\thispagestyle{empty}
\begin{center}
    \Large{Ficha Resumen \\ Trabajo de Grado Maestría en Ingeniería de Software}
\end{center}

\textbf{TÍTULO:} Mecanismo de detección automática de deriva arquitectónica mediante el uso de \textit{Fitness Functions} y análisis estático de grafos de dependencia: un enfoque de gobernanza continua para PyMES de software
\begin{enumerate}
    \item Énfasis: Ingeniería de Software
    \item Área de trabajo: Arquitectura de software, análisis estático de código y gobernanza técnica
    \item Tipo de proyecto (Aplicado, Innovación, Investigación): Investigación aplicada
    \item Estudiante: Juan José Medicis Zambrano
    % TODO ENTREGA: completar con el correo institucional.
    \item Correo electrónico: [COMPLETAR: correo institucional del estudiante]
    % TODO ENTREGA: completar con la dirección y el teléfono de contacto.
    \item Dirección y teléfono: [COMPLETAR: dirección y teléfono del estudiante]
    \item Director: MsC Juan Camilo Parra Martinez
    % TODO ENTREGA: confirmar el tipo de vinculación del director.
    \item Vinculación del director: [COMPLETAR: Planta / Cátedra / Externo]
    % TODO ENTREGA: completar con el correo institucional del director.
    \item Correo electrónico del director: [COMPLETAR: correo institucional del director]
    \item Co-Director (Si aplica): No aplica
    \item Grupo o empresa que lo avala (Si aplica): No aplica
    \item Otros grupos o empresas: No aplica
    \item Palabras clave (al menos 5): deriva arquitectónica; análisis estático; grafos de dependencia; \textit{Fitness Functions}; gobernanza arquitectónica; conformidad estructural
    \item ODS que aplica al proyecto (Agenda 2030): ODS 9 (Industria, innovación e infraestructura)
    \item Fecha de inicio: Mayo de 2026
    \item Resumen: Este proyecto aborda la deriva arquitectónica, entendida como la divergencia progresiva entre la arquitectura prevista de un sistema de software y su implementación real, un fenómeno que incrementa el costo de mantenimiento y erosiona la calidad del producto. Su importancia radica en que las pequeñas y medianas empresas de software, que concentran buena parte del desarrollo en la industria, rara vez disponen de documentación arquitectónica actualizada ni de los recursos para adoptar herramientas de verificación sofisticadas, por lo que acumulan deuda técnica de manera silenciosa. Ante esta problemática, se propone diseñar, construir y evaluar un mecanismo automatizado que detecta derivas arquitectónicas combinando el análisis estático de grafos de dependencia con la ejecución de \textit{Fitness Functions}, definidas como verificaciones objetivas que cuantifican qué tan alineada está la implementación con las reglas estructurales declaradas. El objetivo general consiste en desarrollar y validar experimentalmente dicho mecanismo para clasificar el estado de conformidad estructural de un proyecto sin depender de documentación arquitectónica externa, mientras que los objetivos específicos abarcan la formulación de un catálogo de \textit{Fitness Functions}, la implementación de un componente de análisis estático, la construcción de un motor de validación basado en una Tasa de Conformidad y la ejecución de un experimento controlado sobre un caso de estudio. La evaluación se realizó sobre cinco escenarios de un mismo sistema: una implementación de referencia y cuatro versiones alteradas, tres de ellas con una única deriva inducida para medir la sensibilidad de cada regla por separado. El mecanismo detectó la totalidad de las derivas y no produjo ningún hallazgo sobre la implementación de referencia, superando el criterio de éxito establecido; como complemento, detectó también las 22 derivas inyectadas en cuatro bibliotecas Java reales, sin hallazgos espurios. Los escenarios aislados revelaron, además, que la Tasa de Conformidad diluye las violaciones localizadas por ser una proporción no ponderada, lo que motivó incorporar la severidad del peor hallazgo como criterio complementario de clasificación. Entre las lecciones aprendidas destacan la viabilidad del enfoque con recursos limitados, la importancia de aislar las variables en el diseño experimental —fue esa decisión la que reveló la limitación de la métrica— y la necesidad de calibrar los umbrales de conformidad según el contexto de cada organización.
\end{enumerate}
